%% kapa: exam=1 round=19 year=2008 period=2 form=A subject=economics subject_ko=경제원론 seq=14 number=14
%% source: Q-Net 감정평가사 기출문제 (cst00309) artlSeq=1001962 file=extracted/1008/1001962/2008년 제19회 감정평가사 1차시험 기출문제/2교시.hwp
%% conversion: kordoc 4.24.51 -> tools/gen1.py (auto); review=proofread-sonnet; flags=none
\begin{question}{14}
저소득 계층은 소득 중 음식료품에 대한 지출 비중이 높은 반면에, 고소득 계층은 음식료품에 대한 지출 비중이 상대적으로 낮다고 하자. 이 경우 음식료품에 대한 물품세(sales tax)는?
\begin{choices}
  \item 누진적(progressive)이다.
  \item 역진적(regressive)이다.
  \item 중립적(neutral)이다.
  \item 정액세(lump sum tax)이다.
  \item 부유세(wealth tax)이다.
\end{choices}
\end{question}
\answer{2}
